% !TeX root = ../main.tex

\chapter{Electrons in a Magnetic Field}

\begin{multline*}
  \text{Free electrons in a magnetic field} \\ \longrightarrow
  \text{Bloch electrons (electrons in a periodic potential field) in a
    magnetic field}
\end{multline*}
For metals, the electrons near the Fermi surface can be considered as the
quasifree electrons, where $m^* \to m$.
So, we start from the free electrons first.

\section{Bohr-Sommerfeld Quantization}

\subsection{Classical Hall effect}

In the polar coordinate
\begin{align*}
  \bm r & = r\cos(\omega t)\bm i + r\sin(\omega t)\bm j\\
  \bm{\dot r} & = -r\omega\sin(\omega t)\bm i + r\omega\cos(\omega t)\bm j\\
  \bm{\ddot r} & = -r\omega^2\cos(\omega t)\bm i - r\omega^2\sin(\omega t)\bm j
\end{align*}
where the equation of motion
\[
  m^*\bm{\ddot r} = e\bm{\dot r}\times\bm B
\]
If the magnetic field $B = (0, 0, B)$, then
\[
  m^*(-r\omega^2\cos(\omega t)) = -e\cdot r\omega\cos(\omega t) B
\]
where $r = v/\omega \propto \frac1B$, and the orbit is quantized.

The current density
\[
  \bm j = \frac{ne^2\tau}{m^*}\bm E
\]
while in Drude model, we have
\[
  \bm j = \hat \sigma \bm E
\]
where $\sigma$ is a tensor mentioned in Chapter 1.
If we apply the magnetic field, we use
\[
  \bm E = \frac{m^*}{ne^2\tau}\bm j + \frac1{ne}\bm j \times \bm B
\]
In the Drude model, the diagonal elements are independent from the magnetic
field, while the $xy$ term is linear to the magnetic field.

\subsection{Bohr-Sommerfeld Quantization rule}

The number of wavelength along the trajectory must be integer
\begin{equation}
  \frac1\hbar \oint \bm p \cdot \d\bm r = 2\pi N
\end{equation}
Apply the magnetic field, we need do the replication $\bm p \to \bm p - e\bm A$,
hence
\begin{equation}
  \frac1\hbar \oint \bm p \cdot \d\bm r - \frac1\hbar \oint e\bm A \cdot \d\bm r
= 2\pi N
\end{equation}
where $\bm p = \hbar\bm k$. Using the motion equation becomes
\begin{equation}
  \hbar\bm k = -e\bm r \times \bm B
\end{equation}
The first term of the integral now becomes
\[
  \frac1\hbar \oint \hbar \bm k \cdot \d\bm r
= -\frac e\hbar \oint(\bm r \times \bm B) \cdot \d r
= -\frac e\hbar \oint \bm B \cdot(\d\bm r \times \bm r) = -\frac{2e}{\hbar} \oint \bm B \cdot \d\bm S = \frac{2e}{\Phi}
\]
and the seconde term
\[
  \frac e\hbar \oint \bm A \cdot \d \bm r \xlongequal{\text{Stock' Equation}}
  \frac e\hbar \oint (\nabla \times \bm B) \cdot \d\bm r
= \frac e\hbar \iint \bm B \cdot \d S = \frac e\hbar \Phi
\]
Hence, the quantization condition becomes
\begin{equation}
  \frac e\hbar \Phi = 2\pi N, \quad \Phi = \frac{2\pi\hbar}{e}N
= \frac he N = \phi_0N
\end{equation}
Since $\phi_0 = BS$, so, the area $S$ should be quantized.
\emph{If the magnetic field is changed, the radius of motion will change
  discretely.}

\subsection{Onsager quantization condition}

Going to $k$-space from $r$-space for $\bm B = (0,0,B_z)$
\[
  \hbar|\bm k| = e|\bm r| B, \quad |\bm k| = \frac{|\bm r|}{l_B^2},
\]
where $l_B = \sqrt{\frac\hbar{eB}}$ denotes the magnetic length.
Then, the area
\[
  S_r = \pi r^2, \quad S_k = \pi k^2, \to S_r = l_B^4 S_k
\]
where $S_k$ denotes the area of \emph{Fermi Surface}. Then, the flux
\[
  \phi = N\phi_0, \quad B l_B^4 S_k = N\frac he
\]
Then, we obtain the the quantization condition in $k$-space
\begin{equation}
  l_B^2 S_k = 2\pi N
\end{equation}
To generalized onsager relation
\[
  l_B^2 S_k = 2\pi\ab(N + \frac12 - \frac\gamma{2\pi})
\]
Consider a change in $B$
\[
  \frac\hbar{eB_N} \cdot S_k = 2\pi N, \quad
  \frac\hbar{eB_{N+1}} \cdot S_k = 2\pi (N + 1)
\]
Hence, the flux area change in $k$-space
\begin{equation}
  S_k \cdot \Delta\ab(\frac1B) = \frac{2\pi}{\hbar} e
\end{equation}

\subsection{Quantization of energy}

Starting from
\begin{gather}
  l_B^2 S_k = 2\pi \ab(N + \frac12 - \frac\gamma{2\pi})\\
  S_k = \pi k^2
\end{gather}
The energy $E = \frac{\hbar^2k^2}{2m}$
\[
  \frac\hbar{eB} \cdot \pi \cdot \frac{2mE}{\hbar^2} = 2\pi(N + \frac12), \quad
  E = \hbar\omega\ab(N + \frac12)
\]

\subsection{Full quantum mechanical solution}

The Hamiltonian for a free electron
\[
  \hat H = \frac{\hat p^2}{2m} \xlongrightarrow{\bm p\to\bm p - e\bm A}
  \frac{(\bm p - e\bm A)^2}{2m}
\]
where $\bm B = \nabla \times \bm A$ and use Landau gauge $\bm A = (0,Bx,0)$.
Since $p_y$, $p_z$ commutes with the Hamiltonian, it can be replaced with its
eigenvalue $\hbar k_y$, $\hbar k_z$. The Schr\"odinger equation becomes
\[
  \ab[\frac{p_x^2}{2m^*}
    + \frac1{2m^*}(\hbar k_y - eB_x)^2 + \frac{\hbar^2k^2}{2m^*}]\psi
= E\psi
\]
Rearrange it
\[
  \ab[\frac{p_x^2}{2m^*}
    + \frac1{2m^*}(\hbar k_y - eB_x)^2]\psi
= \ab(E - \frac{\hbar^2k^2}{2m^*})\psi = E'\psi
\]
Denote $x_0 = \frac{\hbar k_y}{eB}$, the Schr\"odinger equation becomes
\begin{equation}
  \ab[\frac{p_x^2}{2m^*}
    + \frac12m^*\omega_c^2(x - x_0)^2]\psi = E'\psi
\end{equation}
The quantized energy
\begin{equation}
  E = \frac{\hbar^2k_z^2}{2m^*} + \ab(N + \frac12)\hbar\omega_c
\end{equation}

Due to the limitation of boundary,
\[
  0 < x_0 = \frac{\hbar k_y}{eB} < L_x
\]
where $k_y = \frac{2\pi}{L_y}n_y$, $\Delta k_y = \frac{2\pi}{L_y}$, it will be
limited within
\[
  0 < k_y < \frac{eBL_x}{\hbar`'}
\]
So, the total number of state
\[
  n = \frac{eBL_x/\hbar}{2\pi/L_y} = \frac{eB(L_xL_y)}{h}
\]
and the density of states
\[
  \text{DOS} = \frac{eB(L_xL_y)/\hbar}{L_xL_y}g_s = \frac{eB}{h}g_s
= \phi_0 B g_s \propto B
\]
Usually, $g_s = 2$, but for like graphene, $g_s > 2$.
By increasing the magnetic field, one can confine all of the electrons within
the same energy (such as the ground) state.

\section{Landau Quantization}

\subsection{Filling factor}

With a given $E_f$, the number of Landau levels filled with states $\nu$.
\begin{equation}
  \nu = \frac{n}{eB/h} = \frac{hn_B}{eB}
\end{equation}
The Landau level will broaden due to scattering.
For a given scattering time $\tau_q$, the broaden energy
\[
  \Delta E \sim \frac\hbar{\tau_q}
\]
due to Heisenberg's uncertainty principle.
When $\Delta E > \hbar\omega_c$, or $\frac\hbar{\tau_c} > \hbar\omega_c$,
$\omega_c\tau_c < 1$. We denotte $\mu = e\tau/m$, then
\[
  \omega_c \cdot \frac{\mu m^*}{e} = \frac{eB}{m^*} \frac{\mu m^*}{e} < 1,
  \quad B \mu_q < 1
\]
So, when $B\mu_q \geq 1$, we could not distinguish the levels.
where $\mu_q$ denotes the quantum mobility. E.g, for
$B = \qty 1\T$, $\mu_q \to \qty{e3}{\cm/\N\cdot\s}$.

\section{Shubnikov-de Hass Oscillators}

\subsection{Zeeman splitting}

Each Landau level will separate into two levels.
\[
  E_c = \hbar \frac{eB}{m^*}, \quad E_z = g\mu_B B
\]
In case that $E_c < E_z$, then $\hbar eB/m^* < g\mu_B B$,
$\hbar eB < gm^*\mu_B B$.

If tune $B$ to $B_z$ and $B_\parallel$, then
\[
  \frac{E_c}{E_B} = \frac{\frac{\hbar eB}{m^*}\cos\theta}{g\mu_B B}
= \frac{\hbar e\cos\theta}{gm^*\mu_B}
\]

\section{Integer Quantum Hall Effect}

\subsection{Edge modes}

The Hamiltonian with the potention $\phi(x)$
\begin{equation}
  \hat H = \frac{p_x^2}{2m^*} + \frac1{2m^*}(\hbar k_y - eBx)^2 + \phi(x)
\end{equation}
The eigenenergy will be
\[
  E(\hbar k) = \ab(N + \frac12)\hbar\omega_c + \phi\ab(\frac{\hbar k_y}{eB})
\]
The current
\[
  I = -e \int_{-\infty}^\infty \frac{\d k}{2\pi} v_y(k)
= \frac{e}{2\pi l_B^2} \int_0^x \d x \frac1{eB} \fdv\phi x
= \frac e{2\pi\hbar} \Delta\mu
\]
where $\pdv E{k_y} = \pdv\phi x \cdot \ab(\frac\hbar{eB})$.
In the bulk, the electrons motion counterclockwise; But on the both edges
of top and top, the electrons will be ``perfectly scattered'' (reflected) by the
edges and do hafl-circle counterclockwise motion.
Obviously, the direction of the currents is different.

\section{Fractional Quantum Hall Effect}

\subsection{Composite fermions}

\[
  B = n_e\phi_e (2j + 1/\nu^*)
\]
Then, $\nu = \frac{\nu^*}{2j\nu^* + 1}$.

If consider $j = 1$, which means that each electrons attached to flux, then
\[
  \nu = \frac13, \rightarrow \nu^* = 1.
\]

\subsection{Wigner crystal}

The electron-electron interaction is so large that forms the lattice.
The bonds are Columb interactions.
When a strong filed is applied on a level, the electrons will be fixed and form
the lattice.